% =============================================================================
% Section 1 — Introduction
% =============================================================================

\section{Introduction}
\label{sec:introduction}

Motor neurone disease (MND), also known as amyotrophic lateral sclerosis (ALS), is a
progressive neurodegenerative condition that destroys upper and lower motor neurons,
resulting in loss of voluntary movement, speech, and ultimately respiratory function.
Global incidence is approximately 2--3 per 100{,}000 person-years, affecting an
estimated 41{,}000 individuals in Australia and 30{,}000 in Brazil~\cite{mehta2018mnd}.
Median survival after diagnosis is 2--5 years; respiratory failure, typically managed
with non-invasive ventilation (NIV), is the leading cause of death. Notably, NIV has
been shown to extend survival by a~factor of approximately three compared with best
pharmacological care alone~\cite{berlowitz2021niv}, making the period of NIV use
clinically significant and prolonged. During this phase, patients remain cognitively
intact while becoming progressively unable to speak or gesture in ways that standard
communication technologies can detect.

Recent work has demonstrated that patient-ventilator asynchrony (PVA) — a mismatch between ventilator cycles and patient breathing effort — is a significant contributor to NIV non-adherence in MND, causing discomfort, sleep fragmentation, and potentially adverse clinical outcomes~\cite{atkeson2011pva, yang2026pva, berlowitz2025tla}. Patients experiencing PVA-related discomfort during nocturnal NIV use cannot readily signal distress to caregivers using existing AAC technologies, which require network connectivity or dedicated hardware unsuited to the home care environment. NeuroAssist AI addresses this gap directly: by providing a local, always-available signalling channel that requires only a commodity webcam, it enables patients to alert caregivers during unassisted periods — including while on NIV at night.

In advanced stages, a~substantial proportion of patients experience a~locked-in-like
condition in which voluntary control is limited to subtle residual facial movements:
small eyelid contractions, brow raises, or microexpressions that persist after limb and
orofacial motor function is lost. These residual gestures represent the last available
voluntary channel for communication. Without a~reliable means to detect and act on them,
patients cannot signal distress, discomfort, or needs to caregivers — with direct
consequences for quality of life and patient safety during unassisted periods.

Existing augmentative and alternative communication (AAC) systems do not adequately
address this population. High-end eye-tracking devices (e.g.\ Tobii Dynavox) and
dedicated AAC software (e.g.\ Grid~3) deliver high accuracy but require stable wireless
connectivity, dedicated hardware costing upwards of USD~3{,}000, and factory-calibrated
setups that cannot adapt to the degrading motor profile of an individual
patient~\cite{beukelman2013aac}. Brain-computer interfaces (BCI) via EEG offer an
alternative, but involve invasive or cumbersome hardware inappropriate for home
deployment~\cite{zander2011passive}. Low-cost camera-based approaches exist in
research~\cite{soukupova2016}, but rely on simple per-frame threshold comparisons that
generate unacceptably high false-positive rates when applied to the involuntary facial
activity present in MND patients during NIV.

We present \textbf{NeuroAssist AI}, a~privacy-preserving, offline-first assistive
communication platform designed for home deployment with a~commodity webcam. The system
transforms residual facial gestures into structured emergency requests using a~locally
executed WebAssembly computer vision pipeline, a~deliberate-gesture state machine, and
a~patient-calibratable detection profile — without transmitting any video off-device.

This paper makes three original contributions:
\begin{enumerate}
  \item \textbf{An offline-first, privacy-preserving detection pipeline.}
    The MediaPipe FaceLandmarker model and its WebAssembly runtime are served from
    local assets; no video frame or biometric data leaves the patient's device.
    The CPU delegate ensures compatibility with commodity hardware regardless of GPU
    availability.

  \item \textbf{A six-state deliberate gesture state machine with patient-specific
    calibration.}
    Rather than a~single EAR threshold, the state machine requires minimum closure
    duration ($\geq$3 frames, ${\approx}$100\,ms) in each of two sequential blinks,
    confirmed re-opening between them, and a~configurable time window. Thresholds are
    stored in a~versioned \texttt{DetectionProfile} and adapted per patient. This design
    specifically targets the false-positive sources relevant to MND: involuntary blinks,
    partial ptosis artefacts, and slow threshold drift.

  \item \textbf{An extensible hexagonal microservices architecture with plug-in gesture
    types.}
    The \texttt{engineFeatureRef} field in the detection profile decouples the gesture
    recognition engine from the command catalogue, allowing new gestures (brow raise,
    single-eye blink) to be registered without modifying existing domain code.
    ArchUnit rules enforced at compile time guarantee that this separation persists as
    the system evolves.
\end{enumerate}

The remainder of the paper is organised as follows. Section~\ref{sec:related} reviews
related work in AAC, facial gesture detection, and edge computing for health.
Section~\ref{sec:architecture} describes the system architecture in detail.
Section~\ref{sec:evaluation} reports the experimental evaluation.
Section~\ref{sec:discussion} discusses limitations and future work, and
Section~\ref{sec:conclusion} concludes.